\section{Code scaling and complexity}
\label{section:preparation:complexity}


\begin{readerguide}
 The complexity discussion is mainly relevant when we conduct an MPI (weak scaling) assessment. 
 In many cases, we might be able to skip it initially.
\end{readerguide}

We have to know or have to be told the computational complexity of the underlying algorithms. 
If we have a pure Monte Carlo setup, then doing 100 times more trials should last around 100 times longer.
If we have a parabolic Partial Differential Equation (PDE) solver with explicit time stepping and a regular grid, we have to know that the problem sizes increases by a factor of $2^d$ if we half the mesh spacing. 
However, these algorithms also have to reduce the time step size, so if we measure over a fixed time span, we actually will have four times more time steps and therefore increase the workload by a factor of $2^{d+2}$ once we half the mesh size.
Finally, if we run an Adaptive Mesh Refinement (AMR) code for a hyperbolic setup with explicit time stepping or a Smoothed-Particle Hydrodynamics (SPH) code, we need some kind of output that tells us exactly how many particles or mesh cells we have and what time step sizes we have picked. In this case, time step sizes go down linearly with mesh spacing.
At the end of the day, we are interested in the \emph{throughput} of the system, i.e.~the time we need to update one quantity of interest.
It is this metric that we will start from in the next steps. 



